\textbf{Phần 1. Trắc nghiệm lựa chọn}
\setcounterpageref{ex}{0}

\begin{ex} %[5D1N7]
Hỗn số $1\dfrac{1}{4}$ được đọc là
\choice
{\True Một và một phần tư}
{Một phần tư}
{Mười một phần tư}
{Một chia một phần tư}
\loigiai{
Hỗn số $1\dfrac{1}{4}$ gồm phần nguyên là $1$ và phần phân số là $\dfrac{1}{4}$, đọc là: Một và một phần tư.\\
\textbf{Đáp án của câu hỏi là A}
}
\end{ex}

\begin{ex} %[5D1N7]
Trong hỗn số $12\dfrac{7}{10}$, phần nguyên và phần phân số lần lượt là
\choice
{$7$ và $\dfrac{12}{10}$}
{$12$ và $\dfrac{10}{7}$}
{\True $12$ và $\dfrac{7}{10}$}
{$10$ và $\dfrac{7}{12}$}
\loigiai{
Hỗn số $12\dfrac{7}{10}$ có phần nguyên là $12$ và phần phân số là $\dfrac{7}{10}$.\\
\textbf{Đáp án của câu hỏi là C}
}
\end{ex}

\begin{ex} %[5D1H7]
Hỗn số $8\dfrac{1}{3}$ được viết dưới dạng tổng của phần nguyên và phần phân số là
\choice
{$8 \times \dfrac{1}{3}$}
{$8 - \dfrac{1}{3}$}
{$8 : \dfrac{1}{3}$}
{\True $8 + \dfrac{1}{3}$}
\loigiai{
Theo định nghĩa hỗn số, $8\dfrac{1}{3} = 8 + \dfrac{1}{3}$.\\
\textbf{Đáp án của câu hỏi là D}
}
\end{ex}

\begin{ex} %[5D1H7]
Chuyển hỗn số $2\dfrac{7}{10}$ thành phân số thập phân, ta được kết quả là
\choice
{$\dfrac{14}{10}$}
{\True $\dfrac{27}{10}$}
{$\dfrac{20}{10}$}
{$\dfrac{72}{10}$}
\loigiai{
Ta có: $2\dfrac{7}{10} = \dfrac{2 \times 10 + 7}{10} = \dfrac{27}{10}$.\\
\textbf{Đáp án của câu hỏi là B}
}
\end{ex}

\begin{ex} %[5D1H7]
Chuyển hỗn số $5\dfrac{1}{10}$ thành phân số thập phân ta được
\choice
{$\dfrac{6}{10}$}
{$\dfrac{15}{10}$}
{\True $\dfrac{51}{10}$}
{$\dfrac{50}{10}$}
\loigiai{
Ta có: $5\dfrac{1}{10} = \dfrac{5 \times 10 + 1}{10} = \dfrac{51}{10}$.\\
\textbf{Đáp án của câu hỏi là C}
}
\end{ex}

\begin{ex} %[5D1H7]
Chuyển hỗn số $1\dfrac{9}{100}$ thành phân số thập phân, ta được
\choice
{\True $\dfrac{109}{100}$}
{$\dfrac{19}{100}$}
{$\dfrac{91}{100}$}
{$\dfrac{190}{100}$}
\loigiai{
Ta có: $1\dfrac{9}{100} = \dfrac{1 \times 100 + 9}{100} = \dfrac{109}{100}$.\\
\textbf{Đáp án của câu hỏi là A}
}
\end{ex}

\begin{ex} %[5D1H7]
Chuyển hỗn số $4\dfrac{9}{100}$ thành phân số thập phân, ta được kết quả là
\choice
{$\dfrac{49}{100}$}
{$\dfrac{490}{100}$}
{$\dfrac{400}{100}$}
{\True $\dfrac{409}{100}$}
\loigiai{
Ta có: $4\dfrac{9}{100} = \dfrac{4 \times 100 + 9}{100} = \dfrac{409}{100}$.\\
\textbf{Đáp án của câu hỏi là D}
}
\end{ex}

\begin{ex} %[5D1H7]
Viết phân số thập phân $\dfrac{57}{10}$ dưới dạng hỗn số, ta được
\choice
{$7\dfrac{5}{10}$}
{$57\dfrac{1}{10}$}
{\True $5\dfrac{7}{10}$}
{$5\dfrac{1}{7}$}
\loigiai{
Thực hiện phép chia: $57 : 10 = 5$ dư $7$. Do đó $\dfrac{57}{10} = 5\dfrac{7}{10}$.\\
\textbf{Đáp án của câu hỏi là C}
}
\end{ex}

\begin{ex} %[5D1H7]
Viết phân số thập phân $\dfrac{703}{100}$ dưới dạng hỗn số, ta được
\choice
{$70\dfrac{3}{100}$}
{\True $7\dfrac{3}{100}$}
{$7\dfrac{30}{100}$}
{$73\dfrac{1}{100}$}
\loigiai{
Ta có: $703 : 100 = 7$ dư $3$. Vậy $\dfrac{703}{100} = 7\dfrac{3}{100}$.\\
\textbf{Đáp án của câu hỏi là B}
}
\end{ex}

\begin{ex} %[5D1H7]
Viết phân số $\dfrac{3\,005}{1\,000}$ dưới dạng hỗn số, ta được
\choice
{$30\dfrac{5}{1\,000}$}
{$3\dfrac{50}{1\,000}$}
{$300\dfrac{5}{1\,000}$}
{\True $3\dfrac{5}{1\,000}$}
\loigiai{
Ta có: $3\,005 : 1\,000 = 3$ dư $5$. Do đó $\dfrac{3\,005}{1\,000} = 3\dfrac{5}{1\,000}$.\\
\textbf{Đáp án của câu hỏi là D}
}
\end{ex}

\begin{ex} %[5D1H7]
Trên tia số giữa hai hỗn số $6\dfrac{2}{10}$ và $6\dfrac{4}{10}$ là vạch biểu diễn hỗn số nào?
\choice
{\True $6\dfrac{3}{10}$}
{$6\dfrac{1}{10}$}
{$6\dfrac{5}{10}$}
{$7\dfrac{3}{10}$}
\loigiai{
Trên tia số có các khoảng chia đều $\dfrac{1}{10}$, điểm nằm giữa $6\dfrac{2}{10}$ và $6\dfrac{4}{10}$ là $6\dfrac{3}{10}$.\\
\textbf{Đáp án của câu hỏi là A}
}
\end{ex}

\begin{ex} %[5D1V7]
Cô Dung có $23$ phong kẹo, mỗi phong kẹo có $10$ viên kẹo. Cô chia đều số kẹo đó cho $10$ bạn. Khẳng định nào sau đây là \textbf{sai}?
\choice
{Mỗi bạn nhận được $23$ viên kẹo}
{\True Mỗi bạn nhận được $2\dfrac{3}{10}$ viên kẹo}
{Mỗi bạn nhận được $\dfrac{23}{10}$ phong kẹo}
{Mỗi bạn nhận được $2\dfrac{3}{10}$ phong kẹo}
\loigiai{
Tổng số viên kẹo là: $23 \times 10 = 230$ viên.\\
Mỗi bạn được: $230 : 10 = 23$ viên kẹo (hoặc $23 : 10 = \dfrac{23}{10} = 2\dfrac{3}{10}$ phong kẹo).\\
Khẳng định ``Mỗi bạn nhận được $2\dfrac{3}{10}$ viên kẹo'' là sai về đơn vị.\\
\textbf{Đáp án của câu hỏi là B}
}
\end{ex}

\textbf{Phần 2. Trắc nghiệm đúng sai}
\setcounterpageref{ex}{0}

\begin{ex} %[5D1H7]
Cho các hỗn số: $A = 5\dfrac{1}{10}$, $B = 1\dfrac{9}{100}$ và các phân số thập phân $M = \dfrac{63}{10}$, $N = \dfrac{703}{100}$. Xét tính đúng sai của các khẳng định sau:
\choiceTF
{\True Chuyển hỗn số $A$ thành phân số thập phân ta được $\dfrac{51}{10}$.}
{Chuyển hỗn số $B$ thành phân số thập phân ta được $\dfrac{19}{100}$.}
{\True Phân số $M$ được viết dưới dạng hỗn số là $6\dfrac{3}{10}$.}
{\True Phân số $N$ được viết dưới dạng hỗn số là $7\dfrac{3}{100}$.}
\loigiai{
\begin{itemize}
    \item $A = 5\dfrac{1}{10} = \dfrac{5 \times 10 + 1}{10} = \dfrac{51}{10}$.
    \item $B = 1\dfrac{9}{100} = \dfrac{1 \times 100 + 9}{100} = \dfrac{109}{100} \neq \dfrac{19}{100}$.
    \item $M = \dfrac{63}{10} = 6 + \dfrac{3}{10} = 6\dfrac{3}{10}$.
    \item $N = \dfrac{703}{100} = 7 + \dfrac{3}{100} = 7\dfrac{3}{100}$.
\end{itemize}
\begin{itemize}
    \item a) \textbf{Đúng}.
    \item b) \textbf{Sai}: $B = \dfrac{109}{100}$.
    \item c) \textbf{Đúng}.
    \item d) \textbf{Đúng}.
\end{itemize}
}
\end{ex}

\begin{ex} %[5D1V7]
Cô Dung chia đều $5$ chiếc bánh cho $4$ bạn học sinh. Xét tính đúng sai của các khẳng định sau:
\choiceTF
{\True Số phần bánh mỗi bạn nhận được là $\dfrac{5}{4}$ cái bánh.}
{Viết dưới dạng hỗn số, mỗi bạn nhận được $1\dfrac{1}{5}$ cái bánh.}
{\True Mỗi bạn nhận được $1$ cái bánh nguyên và thêm $\dfrac{1}{4}$ cái bánh nữa.}
{Để mỗi bạn nhận được đủ $2$ cái bánh thì cô Dung cần thêm $4$ cái bánh nữa.}
\loigiai{
\begin{itemize}
    \item Chia $5$ cái bánh cho $4$ bạn: mỗi bạn được $5 : 4 = \dfrac{5}{4}$ cái bánh.
    \item Ta có $\dfrac{5}{4} = 1\dfrac{1}{4} \neq 1\dfrac{1}{5}$.
    \item Hỗn số $1\dfrac{1}{4}$ gồm phần nguyên là $1$ cái bánh và phần phân số là $\dfrac{1}{4}$ cái bánh.
    \item Để $4$ bạn mỗi bạn có $2$ cái bánh thì cần $4 \times 2 = 8$ cái bánh. Số bánh cần thêm là $8 - 5 = 3$ cái bánh ($\neq 4$ cái bánh).
\end{itemize}
\begin{itemize}
    \item a) \textbf{Đúng}.
    \item b) \textbf{Sai}: Mỗi bạn nhận được $1\dfrac{1}{4}$ cái bánh.
    \item c) \textbf{Đúng}.
    \item d) \textbf{Sai}: Chỉ cần thêm $3$ cái bánh.
\end{itemize}
}
\end{ex}

\textbf{Phần 3. Trả lời ngắn}
\setcounterpageref{ex}{0}

\begin{ex} %[5D1H7]
Chuyển hỗn số $3\dfrac{5}{10}$ thành phân số thập phân có mẫu số là $10$. Hỏi tử số của phân số đó là bao nhiêu?
\shortans{35}
\loigiai{
\textbf{Đáp số:} $35$\\
$3\dfrac{5}{10} = \dfrac{3 \times 10 + 5}{10} = \dfrac{35}{10}$. Tử số là $35$.
}
\end{ex}

\begin{ex} %[5D1H7]
Chuyển hỗn số $4\dfrac{9}{100}$ thành phân số thập phân có mẫu số là $100$. Tử số của phân số nhận được bằng bao nhiêu?
\shortans{409}
\loigiai{
\textbf{Đáp số:} $409$\\
$4\dfrac{9}{100} = \dfrac{4 \times 100 + 9}{100} = \dfrac{409}{100}$. Tử số là $409$.
}
\end{ex}

\begin{ex} %[5D1H7]
Viết phân số thập phân $\dfrac{703}{100}$ dưới dạng hỗn số $a\dfrac{b}{100}$. Tính tích $a \times b$.
\shortans{21}
\loigiai{
\textbf{Đáp số:} $21$\\
Ta có $\dfrac{703}{100} = 7\dfrac{3}{100} \Rightarrow a = 7, b = 3 \Rightarrow a \times b = 21$.
}
\end{ex}

\begin{ex} %[5D1H7]
Viết phân số thập phân $\dfrac{3\,005}{1\,000}$ dưới dạng hỗn số $m\dfrac{n}{1\,000}$. Tính giá trị của $m + n$.
\shortans{8}
\loigiai{
\textbf{Đáp số:} $8$\\
$\dfrac{3\,005}{1\,000} = 3\dfrac{5}{1\,000} \Rightarrow m = 3, n = 5 \Rightarrow m + n = 8$.
}
\end{ex}

\begin{ex} %[5D1V7]
Cô Dung có $23$ phong kẹo, mỗi phong có $10$ viên kẹo. Cô chia đều cho $10$ bạn. Hỏi mỗi bạn nhận được bao nhiêu viên kẹo?
\shortans{23}
\loigiai{
\textbf{Đáp số:} $23$\\
Tổng số kẹo là $23 \times 10 = 230$ viên. Mỗi bạn nhận được $230 : 10 = 23$ viên kẹo.
}
\end{ex}

\begin{ex} %[5D1V7]
Có $17$ lít nước chia đều vào $5$ chiếc can giống nhau. Lượng nước trong mỗi can được viết dưới dạng hỗn số là $a\dfrac{b}{c}\text{ (lít)}$ (với $\dfrac{b}{c}$ là phân số tối giản). Tính giá trị của $a + b + c$.
\shortans{10}
\loigiai{
\textbf{Đáp số:} $10$\\
Mỗi can chứa: $17 : 5 = \dfrac{17}{5} = 3\dfrac{2}{5}\text{ (lít)}$.\\
Do đó $a = 3, b = 2, c = 5 \Rightarrow a + b + c = 3 + 2 + 5 = 10$.
}
\end{ex}
